\printedpage{103}
% head: A. Grothendieck : Dérivateurs, Ch. XII, Édition des Universités de Montpellier II et Paris VII

\noindent\textbf{Corollary 4.} (Assuming W(7 bis) (\footnote[104]{Here we need W(7 bis) in strong form, i.e. $i_0\in W \iff i\in W$ (cf. annotation th. 4).}).) \textit{Let $U \hto[i] X$ be an open immersion such that $T=X\backslash U$ is isomorphic to the one-point category $e$. For $i\in W$ to hold, it is necessary and sufficient that $i$ be $W$-aspheric, i.e. that $U/t$ (where $t$ is the unique object of $T$) be $W$-aspheric.}

It is of course enough to prove that
\[
i\in W \Longrightarrow U/t \ W\text{-aspheric.}
\]

\msp{page 150}
Consider the cartesian square
\[
\begin{tikzcd}[row sep=large,column sep=large]
U/t \arrow[r, hook, "i/t"] \arrow[d, "p"'] & X/t \arrow[d, "q"] \\
U \arrow[r, hook, "i"] & X,
\end{tikzcd}
\]
we are going to prove that it is $W$-cocartesian. Assuming W(7 bis), it will follow that $i\in W$ implies $i/t\in W$, and since $X/t$ is $W$-aspheric (having a final object), this means that $U/t$ is $W$-aspheric, which is what we want to prove.

It is enough, by virtue of corollary 2, to prove that the square is cocartesian in Cat, since $i/t$ is an open immersion ($i$ being one), and $X/t\longrightarrow X$ is a local isomorphism, hence $W$-smooth. We see that $(X/t)\backslash(U/t)\simeq e$ (one-point category), moreover the unique element of $(X/t)\backslash(U/t)$ is the final element $t_0 : t\xrightarrow{\mathrm{id}_t} t$ of $X/t$. Thus $H_{t_0}\simeq e_{U/t}$ (final presheaf), and $p_{!}(H_{t_0})\simeq H'_t \overset{\text{def}}{=} (u\longmapsto \mathrm{Hom}_X(u,t))$, which characterizes the amalgamated sum in Cat
\dots

Corollary 4 is none other than the unproved corollary of proposition 7 (p. 85).

Here is a variant of corollary 4.

\noindent\textbf{Corollary 5.} (Assuming W(7 bis) (\footnote[105]{strong form.}).) \textit{Let $i : U \hto X$ be an open immersion, let $t_0\in\mathrm{Ob}\,T$ be a \textup{closed} point of $T=X\backslash U$, such that $\{t_0\}$}
\msp{page 151}
\textit{is isomorphic to the one-point category, i.e. $\mathrm{Hom}_X(t_0,t_0)=\{\mathrm{id}\}$. Suppose that for every $t\in T\backslash\{t_0\}$, $U/t$ is $W$-aspheric. For $i\in W$ to hold, it is necessary and sufficient that $i$ be aspheric, i.e. that $U/t_0$ also be $W$-aspheric.}

We must prove that $i\in W$ implies that $U/t_0$ is $W$-aspheric. Consider the inclusion
\[
i_0 : U\longrightarrow U' \overset{\text{def}}{=} X\backslash\{t_0\},
\]
the assumption on the $U/t$, for $t\in T\backslash\{t_0\}$, means that $i_0$ is $W$-aspheric. Thus $i_0$ is in $W$, and since the composite $i : U\xrightarrow{i_0} U'\xrightarrow{i'} X$ is in $W$, this implies that
% foot: Version 7 décembre 2002
\printedpage{104}
% head: A. Grothendieck : Dérivateurs, Ch. XII, Édition des Universités de Montpellier II et Paris VII
$i' : U'\longrightarrow X$ is in $W$. We are now in the situation of corollary 4, which implies that $i'$ is $W$-aspheric. Thus $i=i'i_0$ is $W$-aspheric, as a composite of two $W$-aspheric morphisms, q.e.d.

\noindent\textsc{Proof of th. 4.}

\noindent\circled{a} Smoothness criterion for $\varphi_Q : \mathcal{X} = \intdiag{X_0}{X_1}{X_2} \longrightarrow Y$.

Since $\{1\}$, $\{2\}$ are opens of $\Phi = \begin{tikzcd}[row sep=small,column sep=small,baseline=(\tikzcdmatrixname-2-1.base)]
 & 1 \arrow[dl] \\
 0 & \\
 & 2 \arrow[ul]
\end{tikzcd}$, $i_1 : X_1\longrightarrow\mathcal{X}$ and $i_2 : X_2\longrightarrow\mathcal{X}$ are open immersions. Since $\varphi_Q\circ i_1 = f_1$,
\msp{page 152}
$\varphi_Q\circ i_2=f_2$, we already see that
\[
f_1 \text{ and } f_2 \text{ smooth} \iff \varphi_Q \text{ smooth on the open union of } i_1(X_1) \text{ and } i_2(X_2)\text{, i.e. on } \mathcal{X}\backslash i_0(X_0).
\]
It remains to express the smoothness of $\varphi_Q$ at the points $x_0$ of $i_0(X_0)$. Let $x_1$, $x_2$, $y$ be the images of $x_0$ in $X_1$, $X_2$, $Y$ respectively. We must express that the fibers of
\[
(\varphi_Q)/x_0 : \mathcal{X}/x_0\longrightarrow Y/y \tag*{$(*)$}
\]
are aspheric. For this we are going to make $\mathcal{X}/x_0$ and the arrow $(*)$ explicit in terms of the commutative square
\[
(**)\qquad\qquad
\begin{tikzcd}[row sep=large,column sep=large]
X_0/x_0 \arrow[r, "\alpha_1/x_0"] \arrow[d, "\alpha_2/x_0"'] & X_1/x_1 \arrow[d, "f_1/x_1"] \\
X_2/x_2 \arrow[r, "f_2/x_2"] & Y/y,
\end{tikzcd}
\]
induced by $Q$. I say

\noindent\textbf{Lemma 1.} \textit{Let there be a commutative square}
\[
Q \quad = \quad
\begin{tikzcd}[row sep=large,column sep=large]
X_0 \arrow[r, "\alpha_1"] \arrow[d, "\alpha_2"'] & X_1 \arrow[d, "f_1"] \\
X_2 \arrow[r, "f_2"] & Y
\end{tikzcd}
\]
\textit{in} Cat, \textit{$x_0\in\mathrm{Ob}\,X_0$, $x_1$, $x_2$, $y$ its images in $X_1$, $X_2$, $Y$ respectively, $Q/x_0$ the square $(**)$ above, induced by $Q$ on the localized categories $X_0/x_0$ etc., $\mathcal{X} = \intdiag{X_0}{X_1}{X_2}$. We then have a}
% foot: Version 7 décembre 2002
\printedpage{105}
% head: A. Grothendieck : Dérivateurs, Ch. XII, Édition des Universités de Montpellier II et Paris VII
\msp{page 153}
\textit{canonical isomorphism}
\[
\mathcal{X}/x_0 \quad\simeq\quad \intdiag{X_0/x_0}{X_1/x_1}{X_2/x_2} ,
\]
\textit{and the arrow $(\varphi_Q)/x_0$, the localized arrow of $\varphi_Q : \mathcal{X}\longrightarrow Y$ at $x_0$, is identified with the arrow $\varphi_{(Q/x_0)}$ :}
\[
\varphi_{(Q/x_0)} : \quad \intdiag{X_0/x_0}{X_1/x_1}{X_2/x_2} \qquad\longrightarrow\qquad Y/y.
\]

\noindent\textsc{Proof} of the Lemma. We know that
\[
\mathcal{X}/x_0\longrightarrow\mathcal{X}
\]
is Cat-fibered (split) (with discrete fibers), and since
\[
\mathcal{X}\longrightarrow\Phi
\]
is also Cat-fibered (split), it follows that the same is true of the composite arrow
\[
\mathcal{X}/x_0\longrightarrow\Phi.
\]
Thus to ``calculate'' $\mathcal{X}/x_0$, as a category over $\Phi$, we must make explicit its three fibers, and the two base change functors which connect them. We find at once that these fibers are canonically isomorphic to
\[
X_0/x_0, \qquad X_1/x_1, \qquad X_2/x_2,
\]
and we shall admit (!) that the base change functors
\msp{page 154}
\[
X_0/x_0\longrightarrow X_1/x_1, \qquad X_0/x_0\longrightarrow X_2/x_2
\]
(defined up to unique isomorphism) are none other than the localized arrows $\alpha_1/x_0$ and $\alpha_2/x_0$ of $\alpha_1 : X_0\longrightarrow X_1$ and $\alpha_2 : X_0\longrightarrow X_2$. This therefore shows that we have an isomorphism
\[
\mathcal{X}/x_0 \quad\simeq\quad \intdiag{X_0/x_0}{X_1/x_1}{X_2/x_2} ,
\]
% foot: Version 7 décembre 2002
\printedpage{106}
% head: A. Grothendieck : Dérivateurs, Ch. XII, Édition des Universités de Montpellier II et Paris VII
and we shall also admit that the arrow $(\varphi_Q)/x_0 : \mathcal{X}/x_0\longrightarrow Y/y$ is identified with $\varphi_{(Q/x_0)}$.

To express the $W$-smoothness of $\varphi_Q$ at $x_0\in\mathrm{Ob}\,X_0\subseteq\mathrm{Ob}\,\mathcal{X}$, we are thus reduced to expressing the $W$-asphericity of the fibers of
\[
\varphi_{(Q/x_0)} : \quad \intdiag{X_0/x_0}{X_1/x_1}{X_2/x_2} \qquad\longrightarrow\qquad Y/y.
\]
To make these fibers explicit, we need another auxiliary lemma

\noindent\textbf{Lemma 2} (\footnote[106]{NB This is b) of th. 4.}). \textit{Let there be a commutative square in} Cat
\[
Q \quad = \quad
\left(
\begin{tikzcd}[row sep=large,column sep=large]
X_0 \arrow[r, "\alpha_1"] \arrow[d, "\alpha_2"'] & X_1 \arrow[d, "f_1"] \\
X_2 \arrow[r, "f_2"] & Y
\end{tikzcd}
\right),
\]
\textit{whence}
\[
\varphi_Q : \quad \left( \intdiag{X_0}{X_1}{X_2} \right) \qquad\longrightarrow\qquad Y,
\]

\msp{page 155}
\textit{and let $y\in\mathrm{Ob}\,Y$. Then the fiber of $\varphi_Q$ at $y$}
\[
\mathcal{X}_y = \varphi_Q^{-1}(\ \underbrace{\{y\}}_{\text{one-point category}}\ )
\]
\textit{is given by (as a category over $\Phi$)}
\[
\mathcal{X}_y \quad\simeq\quad \intdiag{X_{0\,y}}{X_{1\,y}}{X_{2\,y}} .
\]

\noindent\textsc{Proof.} It is immediate that we have
\[
\mathcal{X}_y\cap\underbrace{\mathcal{X}_s}_{X_s}\simeq(X_s)_y \qquad\quad \text{for } s\in\mathrm{Ob}\,\Phi=\{0,1,2\},
\]
% foot: Version 7 décembre 2002
\printedpage{107}
% head: A. Grothendieck : Dérivateurs, Ch. XII, Édition des Universités de Montpellier II et Paris VII
and (chance does things well) the base change functors
\[
X_0\longrightarrow X_1, \qquad X_0\longrightarrow X_2
\]
send the (nonfull) subcategory $X_{0\,y}$ of $X_0$ into the (nonfull) subcategories $X_{1\,y}$ and $X_{2\,y}$ of $X_1$ and $X_2$ respectively.
\[
\begin{tikzcd}[row sep=large,column sep=large]
X_{s\,y} \arrow[r] \arrow[d] \arrow[dr, phantom, "\scriptstyle\text{cart.}"] & \mathcal{X}_y \arrow[r] \arrow[d] \arrow[dr, phantom, "\scriptstyle\text{cart.}"] & \{y\} \arrow[d] \\
X_s \arrow[r] \arrow[d] \arrow[dr, phantom, "\scriptstyle\text{cart.}"] & \mathcal{X} \arrow[r] \arrow[d] & Y \\
\{s\} \arrow[r] & \Phi &
\end{tikzcd}
\]
We thus have a fibered subcategory of $\mathcal{X}$, subject to checking that the strictly cartesian arrows
\[
x_0\longrightarrow\alpha_1(x_0), \qquad x_0\longrightarrow\alpha_2(x_0)
\]
in $\mathcal{X}$, for $x_0$ in $X_{0\,y}$, are in $\mathcal{X}_y$, i.e. are over the identity arrow of $y$. Now that is the essential property
\msp{page 156}
of the functor $\varphi_Q : \mathcal{X}\longrightarrow Y$, which we forgot to make explicit in 3.2. It really seems to work (like clockwork!)

Let us apply Lemma 2 to the square diagram $Q/x_0$ and to $\varphi_{(Q/x_0)}$ (page 154), we thus see that the fibers to be examined at the $u\in\mathrm{Ob}\,Y/y$ ($u : y'\longrightarrow y$) can be written as the integrals
\[
\mathcal{X}(x_0,u) \quad\overset{\text{def}}{=}\quad \intdiag{(X_0/x_0)_u}{(X_1/x_1)_u}{(X_2/x_2)_u} .
\]
Let us now use the fact that $f_1$, $f_2$ are $W$-smooth (expressing that $\varphi_Q$ is $W$-smooth in $\mathcal{X}\backslash X_0$), hence that $(X_1/x_1)_u$ and $(X_2/x_2)_u$ are $W$-aspheric. In the $W$-cocartesian square
\[
\begin{tikzcd}[row sep=large,column sep=small]
 & (X_1/x_1)_u \arrow[dr, "{i_1(x_0,u)}"] & \\
(X_0/x_0)_u \arrow[ur, "{\alpha_1(x_0,u)}"] \arrow[dr, "{\alpha_2(x_0,u)}"'] & & (\mathcal{X}/x_0)_u\simeq\mathcal{X}(x_0,u) \\
 & (X_2/x_2)_u \arrow[ur, "{i_2(x_0,u)}"'] &
\end{tikzcd}
\]
% foot: Version 7 décembre 2002